\section{Background and State of the Art}

Neurodegenerative diseases impose a growing societal and healthcare
burden. Blood-based proteomics has emerged as a scalable route to
capture systemic and brain-related molecular variation. At the same
time, recent machine learning advances have enabled high-dimensional
pattern extraction, but major limitations remain in transportability,
biological interpretability, and leakage-free evaluation.

A central gap is the lack of principled frameworks that jointly address:
\begin{itemize}
  \item cross-disease modeling,
  \item cross-platform proteomics harmonization,
  \item uncertainty-aware prediction, and
  \item clinically meaningful validation.
\end{itemize}

\subsection{Knowledge gap}

Existing studies often focus on a single disease, a single cohort, or a
single proteomics platform. This limits generalizability and weakens
downstream biological interpretation.

% Insert with % Insert with % Insert with \input{optional/01-example-figure} inside a subsection.
\begin{figure}[H]
  \centering
  \IfFileExists{figures/example_figure.pdf}{%
    \includegraphics[width=0.82\textwidth]{figures/example_figure.pdf}%
  }{%
    \fbox{%
      \parbox[c][4cm][c]{0.78\textwidth}{\centering
      Placeholder figure: add
      \texttt{figures/example\_figure.pdf}}%
    }%
  }
  \caption{Example inserted figure. Replace the file with your own
  figure in \texttt{figures/}.}
  \label{fig:inserted}
\end{figure}
 inside a subsection.
\begin{figure}[H]
  \centering
  \IfFileExists{figures/example_figure.pdf}{%
    \includegraphics[width=0.82\textwidth]{figures/example_figure.pdf}%
  }{%
    \fbox{%
      \parbox[c][4cm][c]{0.78\textwidth}{\centering
      Placeholder figure: add
      \texttt{figures/example\_figure.pdf}}%
    }%
  }
  \caption{Example inserted figure. Replace the file with your own
  figure in \texttt{figures/}.}
  \label{fig:inserted}
\end{figure}
 inside a subsection.
\begin{figure}[H]
  \centering
  \IfFileExists{figures/example_figure.pdf}{%
    \includegraphics[width=0.82\textwidth]{figures/example_figure.pdf}%
  }{%
    \fbox{%
      \parbox[c][4cm][c]{0.78\textwidth}{\centering
      Placeholder figure: add
      \texttt{figures/example\_figure.pdf}}%
    }%
  }
  \caption{Example inserted figure. Replace the file with your own
  figure in \texttt{figures/}.}
  \label{fig:inserted}
\end{figure}

% Insert with % Insert with % Insert with \input{optional/02-example-tikz} inside a subsection.
\begin{figure}[H]
\centering
\begin{tikzpicture}[
  node distance=1.8cm and 1.8cm,
  box/.style={
    draw,
    rounded corners,
    align=center,
    minimum width=3.0cm,
    minimum height=0.9cm
  },
  >=Latex
]
\node[box] (data) {Cohorts \\ Proteomics + Clinical Data};
\node[box, right=of data] (model) {Representation \\ Learning Model};
\node[box, right=of model] (outcome)
  {Diagnosis / Prognosis \\ Biological Interpretation};

\draw[->] (data) -- (model);
\draw[->] (model) -- (outcome);
\end{tikzpicture}
\caption{Example schematic drawn directly in LaTeX using TikZ.}
\label{fig:tikz}
\end{figure}
 inside a subsection.
\begin{figure}[H]
\centering
\begin{tikzpicture}[
  node distance=1.8cm and 1.8cm,
  box/.style={
    draw,
    rounded corners,
    align=center,
    minimum width=3.0cm,
    minimum height=0.9cm
  },
  >=Latex
]
\node[box] (data) {Cohorts \\ Proteomics + Clinical Data};
\node[box, right=of data] (model) {Representation \\ Learning Model};
\node[box, right=of model] (outcome)
  {Diagnosis / Prognosis \\ Biological Interpretation};

\draw[->] (data) -- (model);
\draw[->] (model) -- (outcome);
\end{tikzpicture}
\caption{Example schematic drawn directly in LaTeX using TikZ.}
\label{fig:tikz}
\end{figure}
 inside a subsection.
\begin{figure}[H]
\centering
\begin{tikzpicture}[
  node distance=1.8cm and 1.8cm,
  box/.style={
    draw,
    rounded corners,
    align=center,
    minimum width=3.0cm,
    minimum height=0.9cm
  },
  >=Latex
]
\node[box] (data) {Cohorts \\ Proteomics + Clinical Data};
\node[box, right=of data] (model) {Representation \\ Learning Model};
\node[box, right=of model] (outcome)
  {Diagnosis / Prognosis \\ Biological Interpretation};

\draw[->] (data) -- (model);
\draw[->] (model) -- (outcome);
\end{tikzpicture}
\caption{Example schematic drawn directly in LaTeX using TikZ.}
\label{fig:tikz}
\end{figure}


\subsection{Why this project is needed}

A unified framework could support both biomarker discovery and disease
mechanism mapping while remaining compatible with real-world clinical
translation.
